\documentclass[11pt,a4paper]{article}

% ---------- packages ----------
\usepackage[utf8]{inputenc}
\usepackage[T1]{fontenc}
\usepackage{lmodern}
\usepackage[margin=2.2cm]{geometry}
\usepackage{xcolor}
\usepackage{titlesec}
\usepackage{enumitem}
\usepackage{booktabs}
\usepackage{array}
\usepackage{tabularx}
\usepackage{fancyhdr}
\usepackage{tcolorbox}
\usepackage{hyperref}
\usepackage{graphicx}
\usepackage{parskip}
\tcbuselibrary{skins,breakable}

% ---------- colors ----------
\definecolor{brand}{HTML}{2F6FED}
\definecolor{brandDark}{HTML}{1B3A8A}
\definecolor{ink}{HTML}{1F2933}
\definecolor{muted}{HTML}{52606D}
\definecolor{rule}{HTML}{D9E2EC}
\definecolor{ok}{HTML}{1E8E5A}
\definecolor{warn}{HTML}{C77700}
\definecolor{bad}{HTML}{C0392B}
\definecolor{panel}{HTML}{F4F7FB}

\color{ink}

% ---------- hyperref ----------
\hypersetup{
  colorlinks=true,
  linkcolor=brand,
  urlcolor=brand,
  pdftitle={Life Planner --- Architecture \& Integration Roadmap},
  pdfauthor={Life Planner}
}

% ---------- section styling ----------
\titleformat{\section}
  {\color{brandDark}\Large\bfseries}{\thesection}{0.6em}{}
\titleformat{\subsection}
  {\color{brand}\large\bfseries}{\thesubsection}{0.6em}{}
\titlespacing*{\section}{0pt}{1.4em}{0.5em}

% ---------- header / footer ----------
\pagestyle{fancy}
\fancyhf{}
\renewcommand{\headrulewidth}{0.4pt}
\renewcommand{\footrulewidth}{0pt}
\fancyhead[L]{\small\color{muted}Life Planner}
\fancyhead[R]{\small\color{muted}Architecture \& Integration Roadmap}
\fancyfoot[C]{\small\color{muted}\thepage}

% ---------- helpers ----------
\newcommand{\yes}{\textcolor{ok}{\textbf{Yes}}}
\newcommand{\no}{\textcolor{bad}{\textbf{No}}}
\newcommand{\maybe}{\textcolor{warn}{\textbf{Partial}}}

\newtcolorbox{keybox}[1]{
  breakable, enhanced, colback=panel, colframe=brand,
  boxrule=0.4pt, arc=2pt, left=10pt, right=10pt, top=8pt, bottom=8pt,
  fonttitle=\bfseries\color{white},
  coltitle=white, title={#1},
  attach title to upper={\par\medskip},
  colbacktitle=brand
}

\newtcolorbox{notebox}{
  breakable, enhanced, colback=panel, colframe=rule,
  boxrule=0.4pt, arc=2pt, left=10pt, right=10pt, top=8pt, bottom=8pt
}

\setlist[itemize]{leftmargin=1.4em, itemsep=2pt, topsep=3pt}
\setlist[enumerate]{leftmargin=1.6em, itemsep=3pt, topsep=3pt}

% ============================================================
\begin{document}

% ---------- title ----------
\begin{center}
  {\color{brandDark}\Huge\bfseries Life Planner}\\[0.4em]
  {\color{brand}\LARGE Architecture \& Integration Roadmap}\\[0.8em]
  {\color{muted}\large PWA, VPS Deployment, and a Pluggable Data-Integration Hub}\\[0.6em]
  {\color{muted} Prepared 2026-06-16}
\end{center}

\vspace{0.4em}
\hrule height 1pt
\vspace{1.4em}

% ============================================================
\section{System Overview: Three Services}

The application runs as three cooperating services. Each layer depends on the
one below it.

\begin{keybox}{The three running services}
\begin{enumerate}
  \item \textbf{PostgreSQL} (port \texttt{5432}) --- the database and source of truth.
  \item \textbf{Backend} --- FastAPI via \texttt{uvicorn} (port \texttt{8000}); reads and writes Postgres.
  \item \textbf{Frontend} --- Vite dev server (port \texttt{5173}); the browser UI that calls the backend.
\end{enumerate}
\end{keybox}

The \texttt{start.sh} script orchestrates all three for local development:
it ensures Postgres is up, runs database migrations, launches the backend
(waiting for port 8000), launches the frontend (waiting for port 5173), then
monitors both and shuts everything down cleanly on \texttt{Ctrl+C}.

\begin{notebox}
\textbf{Fixed bug:} \texttt{start.sh} previously ran
\texttt{alembic upgrade head -q}. The \texttt{-q} flag is not valid for Alembic,
so the migration step errored and the script aborted before starting any
servers. Removing \texttt{-q} (matching \texttt{setup.sh}) restored it.
\end{notebox}

% ============================================================
\section{Progressive Web App (PWA)}

A PWA is the existing frontend plus two pieces:

\begin{itemize}
  \item \textbf{Web App Manifest} --- name, icons, theme color, \texttt{display: standalone};
        this makes the app \emph{installable} (Add to Home Screen, fullscreen launch).
  \item \textbf{Service Worker} --- a background script that intercepts network
        requests and serves cached assets; this enables \emph{offline} use and faster loads.
\end{itemize}

With Vite these are generated by the \texttt{vite-plugin-pwa} plugin
(a Workbox wrapper) rather than written by hand.

\subsection{The key caveat: a PWA caches the frontend, not the data}

\begin{center}
\renewcommand{\arraystretch}{1.35}
\begin{tabularx}{\textwidth}{>{\raggedright\arraybackslash}p{4.2cm} >{\raggedright\arraybackslash}X >{\centering\arraybackslash}p{3.2cm}}
\toprule
\textbf{Layer} & \textbf{Lives where} & \textbf{Cacheable?} \\
\midrule
App shell (HTML/JS/CSS, icons) & Static build output & \yes \\
Planner data (tasks, focuses, areas) & Postgres, via \texttt{/api} on :8000 & \maybe \\
\bottomrule
\end{tabularx}
\end{center}

With a basic setup the app loads instantly and offline (the shell is cached),
but any \texttt{/api} call still needs the backend and Postgres. True offline
support means caching API responses in IndexedDB, reading from it when offline,
queuing writes, and syncing on reconnect --- a separate, larger project.

\begin{itemize}
  \item \textbf{Tier 1 (small):} installable app + cached shell + cached GET responses.
        Feels native; still needs the backend for fresh data and writes.
  \item \textbf{Tier 2 (large):} full local-first / offline-write sync via IndexedDB.
\end{itemize}

\textbf{Decision:} start with \textbf{no offline support} --- Tier 1 only.

% ============================================================
\section{VPS Deployment}

Deploying to a VPS resolves the PWA plumbing problems.

\begin{center}
\renewcommand{\arraystretch}{1.3}
\begin{tabularx}{\textwidth}{>{\raggedright\arraybackslash}p{5.5cm} X}
\toprule
\textbf{The VPS solves} & \textbf{The VPS does \emph{not} solve} \\
\midrule
HTTPS via a domain + Let's Encrypt, so service workers register and the PWA is installable anywhere &
Offline data (Tier 2) --- identical problem; a remote backend makes caching \emph{more} valuable but no easier \\
Single origin: a reverse proxy serves \texttt{dist/} and proxies \texttt{/api} on one domain & \\
Public reachability over mobile data, not just home Wi-Fi & \\
\bottomrule
\end{tabularx}
\end{center}

\subsection{Production runtime differs from \texttt{start.sh}}

\texttt{start.sh} is a development runner (\texttt{--reload}, Vite dev server,
crash-monitor loop). Production on a VPS looks like:

\begin{itemize}
  \item \textbf{Frontend:} \texttt{vite build} $\rightarrow$ static \texttt{dist/}, served by Caddy/nginx.
  \item \textbf{Backend:} \texttt{uvicorn}/gunicorn under \textbf{systemd} (auto-restart), no \texttt{--reload}.
  \item \textbf{Postgres:} already a systemd service.
  \item \textbf{Reverse proxy:} Caddy terminates HTTPS, serves static, proxies \texttt{/api}.
\end{itemize}

% ============================================================
\section{Mobile, Health \& AI Integration}

\begin{keybox}{The governing insight}
The backend (FastAPI + Postgres) stays as the \textbf{hub and source of truth}.
Samsung Health and the OS assistant are \textbf{on-device APIs, not cloud APIs} ---
there is no server-to-server health API the VPS can call. A device-native layer
reads the data and pushes it to the backend over HTTP:

\medskip
\centering
\texttt{[native: Health Connect / assistant]} $\rightarrow$ HTTP $\rightarrow$ \texttt{[FastAPI on VPS]} $\rightarrow$ \texttt{[Postgres]}
\end{keybox}

\subsection{Where each integration lands}

\begin{center}
\renewcommand{\arraystretch}{1.35}
\begin{tabularx}{\textwidth}{>{\raggedright\arraybackslash}X >{\centering\arraybackslash}p{2.6cm} >{\raggedright\arraybackslash}p{6.2cm}}
\toprule
\textbf{Goal} & \textbf{From a PWA?} & \textbf{What it really requires} \\
\midrule
Installable app, push notifications & \yes & PWA on the VPS (HTTPS) \\
Read Samsung Health steps/activity & \no & Native Android reading \textbf{Health Connect} \\
Smartwatch (Galaxy Watch) & \no & Native \textbf{Wear OS} app (Kotlin/Compose) \\
``Gemini aware of it \& manages it'' & \maybe & Mostly \emph{your own} in-app AI, not system Gemini \\
\bottomrule
\end{tabularx}
\end{center}

\subsection{Samsung Health / activity goals}
Modern Android health data flows through \textbf{Health Connect} (Google's
on-device store that Samsung Health writes into).
\begin{itemize}
  \item It is a \textbf{device API}, reachable only from a native Android app
        (or a cross-platform wrapper with a plugin) --- not from a web PWA.
  \item Production read access needs a \textbf{Google review / allowlist}
        (privacy policy, declared permissions).
  \item Flow: native app reads steps/workouts $\rightarrow$ \texttt{POST}s to a backend
        endpoint $\rightarrow$ backend matches them to activity-type goals and updates progress.
\end{itemize}

\subsection{Gemini / ``aware and manages it''}
\begin{itemize}
  \item \textbf{System Gemini managing a third-party app:} largely not openly
        available. Limited surface via \emph{App Actions / App Functions}
        (the assistant can \emph{launch} app flows via intents); deep control is gated.
  \item \textbf{Your own in-app assistant} (fully in your control): expose the
        FastAPI endpoints as \textbf{tools / function-calling} to an LLM
        (the Claude API with tool use), letting users chat to create and manage
        goals. The clean backend API is what makes this straightforward.
\end{itemize}

\subsection{Mobile app + smartwatch}
\begin{itemize}
  \item A PWA installs on the phone but cannot reach Health Connect, sensors, or
        Wear OS. Galaxy Watch (4+) runs Wear OS; watch apps are native --- there is
        no ``PWA on a watch.''
  \item Lowest-effort path: \textbf{wrap the existing React frontend in Capacitor}
        $\rightarrow$ native Android shell + a Health Connect plugin, reusing the
        current UI. A small Wear OS companion syncs simple data via the backend.
\end{itemize}

% ============================================================
\section{Data Integrations: A Pluggable Catalog}

The long-term value of this project is concentration: many personal-data sources,
each feeding goals, combined in one place. The differentiator is not any single
integration but the \textbf{hub} --- and a clean pattern for adding more over time.

\begin{keybox}{Two integration patterns}
Every external source falls into one of two categories, and which one decides how
hard it is to build:

\medskip
\textbf{1. Server-side (cloud API).} The VPS calls the provider directly over HTTP.
No phone involved. Often only a \emph{username} (public data) or a one-time OAuth
token. \emph{Easiest to ship.} Example: \textbf{GitHub}.

\medskip
\textbf{2. Device-side (on-device API).} Only a native app on the phone can read it;
it then pushes to the backend. Needs the Capacitor/native layer and OS permissions.
Example: \textbf{Health Connect}, \textbf{screen time}.
\end{keybox}

\subsection{GitHub --- a server-side example (easy win)}
GitHub exposes a full REST/GraphQL API, so the backend can pull contribution data
on a schedule with \textbf{no device involved}.
\begin{itemize}
  \item \textbf{Public data} (contribution graph, public commits, public repos) needs
        only a \textbf{username} --- ideal for a zero-friction ``coding streak'' goal.
  \item \textbf{Private activity} (private commits, full counts) needs a personal
        access token or OAuth.
  \item Flow: a scheduled backend job hits the GitHub API $\rightarrow$ updates
        ``commit / project'' goals (e.g.\ \emph{commit every day}, \emph{N PRs this month}).
\end{itemize}

\subsection{Screen time --- device-side, with a server-side shortcut}
Screen-time goals (e.g.\ \emph{under 5 hours/day}) work as described: set the target
on the website, and a native layer reads usage and reports back.
\begin{itemize}
  \item \textbf{Android (device-side):} \texttt{UsageStatsManager} / Digital Wellbeing,
        gated by the special \texttt{PACKAGE\_USAGE\_STATS} permission the user grants
        in Settings. The Capacitor app reads it and \texttt{POST}s daily totals.
  \item \textbf{iOS:} Screen Time (\texttt{DeviceActivity}/\texttt{FamilyControls}) is
        heavily restricted --- data largely cannot leave the device.
  \item \textbf{Server-side shortcut:} \textbf{RescueTime} offers a cross-platform
        cloud API for productivity/active time --- a way to ship screen-time-style
        goals \emph{before} building the native usage-stats path.
\end{itemize}

\subsection{Catalog of useful data platforms}

\textbf{Server-side (cloud API from the VPS --- ship these first):}
\begin{center}
\renewcommand{\arraystretch}{1.3}
\begin{tabularx}{\textwidth}{>{\raggedright\arraybackslash}p{3.0cm} >{\raggedright\arraybackslash}X >{\raggedright\arraybackslash}p{3.4cm}}
\toprule
\textbf{Platform} & \textbf{Goal type it powers} & \textbf{Auth} \\
\midrule
GitHub / GitLab & Commits, PRs, project streaks & Username / OAuth \\
WakaTime & Coding time per language / project & OAuth \\
Strava & Runs, rides, distance (no native needed) & OAuth \\
Fitbit / Oura / Withings & Steps, sleep, heart rate (cloud route) & OAuth \\
Spotify & Listening, new-artist discovery & OAuth \\
Goodreads / Hardcover / OpenLibrary & Books read, reading pace & Key / username \\
LeetCode & Problems solved, contest streaks & Username (unofficial) \\
RescueTime & Productive vs distracting time & API key \\
Todoist / Notion / Google Calendar & Task \& schedule sync & OAuth \\
Plaid / bank APIs & Budgeting, savings targets & OAuth \\
OpenWeather & Ambient context (enriches goals) & API key \\
\bottomrule
\end{tabularx}
\end{center}

\textbf{Device-side (require the native layer):}
\begin{center}
\renewcommand{\arraystretch}{1.3}
\begin{tabularx}{\textwidth}{>{\raggedright\arraybackslash}p{3.6cm} >{\raggedright\arraybackslash}X >{\raggedright\arraybackslash}p{3.0cm}}
\toprule
\textbf{Source} & \textbf{Goal type it powers} & \textbf{Platform} \\
\midrule
Health Connect & Steps, sleep, heart rate, workouts & Android \\
UsageStatsManager / Digital Wellbeing & Screen time, app usage & Android \\
HealthKit & Health \& fitness metrics & iOS \\
Location / geofencing & Place-based habits, check-ins & Android / iOS \\
\bottomrule
\end{tabularx}
\end{center}

\begin{notebox}
\textbf{Note on health sources:} Strava, Fitbit, Oura, and Withings expose
\emph{cloud} APIs, so activity/health goals can be shipped \textbf{server-side via
OAuth} --- a faster route than the native Health Connect path, which remains the
option for Samsung Health specifically.
\end{notebox}

\subsection{Designing for ``add more later''}
To keep onboarding new sources cheap, treat every integration as a uniform
\emph{provider}:
\begin{itemize}
  \item A \texttt{providers} table (name, kind: \texttt{server}/\texttt{device}, auth
        type) and per-user \texttt{connections} (tokens / usernames).
  \item A normalized \texttt{POST /api/metrics} ingestion shape so both scheduled
        server-side pulls and device-side pushes write the same way.
  \item Goals reference a \texttt{metric} (e.g.\ \texttt{github.commits},
        \texttt{screen.minutes}, \texttt{steps}) rather than a hard-coded source ---
        so a new platform is mostly config, not new plumbing.
\end{itemize}

% ============================================================
\section{Recommended Sequencing}

\begin{enumerate}
  \item \textbf{VPS + HTTPS + reverse proxy} (Caddy serving \texttt{dist/}, proxying \texttt{/api}).
  \item \textbf{Tier-1 PWA} (installable, push notifications) --- falls out of step 1.
  \item \textbf{Server-side integrations} (GitHub first, then WakaTime / Strava / RescueTime)
        --- highest value per effort, no native code.
  \item \textbf{In-app AI assistant} via Claude API function-calling over existing endpoints.
  \item \textbf{Capacitor wrap + Health Connect} --- first real native step; unlocks activity goals.
  \item \textbf{Device-side screen time} (\texttt{UsageStatsManager}) once the native layer exists.
  \item \textbf{Wear OS companion} --- last, narrowest scope.
\end{enumerate}

\subsection{Cheap backend decisions to make now}
\begin{itemize}
  \item \textbf{Model goals with a tracking source} --- e.g. \texttt{manual} vs
        \texttt{activity\_metric} (steps, distance, active minutes), so activity
        goals can later auto-update.
  \item \textbf{Keep an ingestion endpoint in mind} --- e.g. \texttt{POST /api/activity}
        accepting metric samples; the native layer feeds it later.
\end{itemize}

\vspace{1em}
\begin{notebox}
\centering\itshape
The three services and the backend-as-hub model stay constant. Everything new ---
PWA, VPS, native app, watch, assistant, and every data integration --- attaches
around that core as a uniform provider rather than replacing it. Each source added
makes the combined whole more unique.
\end{notebox}

\end{document}
